% =====================================================================
%  Capítulo 3 — contenido
%
%  Sólo el cuerpo del capítulo: sin preámbulo, sin \begin{document},
%  sin portada, sin índices y sin bibliografía (los pone el documento
%  contenedor: cap03.tex para la entrega suelta, libro.tex para el libro).
%
%  Convenciones:
%   * Figuras en capitulos/cap03/figuras/ con prefijo cap03-
%     y citadas sólo por nombre: \includegraphics{cap03-esquema.png}
%   * Etiquetas: fig:cap03:..., tab:cap03:..., sec:cap03:..., lst:cap03:...
%   * Citas con \cite{clave} sobre bib/referencias.bib
% =====================================================================

\chapter{\TituloCapTres}
\label{cap:03}

\section{Introducción}
\label{sec:cap03:introduccion}

Contenido pendiente.

\section{Desarrollo}
\label{sec:cap03:desarrollo}

Contenido pendiente.

% Ejemplo de figura (descomenta y ajusta):
% \Figura{cap03-ejemplo.png}{0.7}{Pie de la figura.}{capjemplo}

% Ejemplo de tabla (descomenta y ajusta):
% \begin{table}[H]
%     \centering
%     \caption{Pie de la tabla.}
%     \label{tab:capjemplo}
%     \begin{tabular}{lrr}
%         \toprule
%         \textbf{Columna} & \textbf{Valor} & \textbf{Unidad} \\
%         \midrule
%         Ejemplo & 0.00 & s \\
%         \bottomrule
%     \end{tabular}
% \end{table}

\section{Conclusiones del capítulo}
\label{sec:cap03onclusiones}

Contenido pendiente.
